% !TEX root = ../main.tex

\chapter{绪论}

\section{研究背景：跨链资产交换}

区块链技术通过分布式账本、密码学签名和共识机制实现了去中心化的价值记录与转移~\cite{nakamoto2008bitcoin,zheng2018blockchain}。随着以太坊、比特币、BSC、Conflux 等公链以及 Hyperledger Fabric 等联盟链的发展，区块链逐渐形成公链与联盟链并存的多链生态格局~\cite{lohachab2021towards}。不同形式和技术的区块链基于其在共识性能、节点组织形式上的不同，适用的应用场景具有显著差异。比如，以Etheruem, Conflux等区块链基于他们强大的合约虚拟环境和庞大的节点支持，在面向以智能合约执行的分布式业务中具有显著优势；而联盟链基于其相比公有链更高的交易吞吐率，侧重高性能交易处理；另外，以门罗币，零币为代表的区块链通过密码手段代替脚本机制，则为以安全性和价值保护为核心目标的业务提供极高的安全保障。
需要强调的是，门罗币（Monero）、Zcash 等无脚本/强隐私区块链同样存在明确的跨链互通诉求，而非孤立于多链生态之外。其核心矛盾在于：隐私币用户既希望保有链上隐私与资产安全，又需要在不依赖中心化托管交易所的前提下，将隐私资产兑换为比特币、以太坊等主流资产，或反向完成入金；若只能通过中心化平台出入金，则会重新引入托管风险并削弱隐私设计初衷。Thyagarajan 等人提出的 Universal Atomic Swaps 指出，Monero、Zcash 屏蔽地址以及 Mimblewimble 等系统因不支持哈希锁与时间锁脚本，传统 HTLC 原子交换无法直接适用，因而需要仅依赖签名验证能力的通用跨链方案~\cite{thyagarajan2022universal}。面向门罗币，Thyagarajan、Malavolta 等人进一步提出 PayMo，基于可验证定时可链接环签名构造与门罗币交易格式兼容的无脚本原子交换，使其可与比特币、以太坊等多种主流币种安全互换~\cite{thyagarajan2022paymo}；Hoenisch 等人提出的 LightSwap 则给出比特币与门罗币之间的轻量原子交换协议，说明该类互通具备可落地的工程场景~\cite{hoenisch2022lightswap}。针对 Zcash，Sanchez 等人提出 Bridging Sapling（ZCLAIM），面向 Sapling 屏蔽池设计隐私跨链资产迁移框架，表明隐私币用户同样需要在尽量保持隐私的同时完成跨链资产流转~\cite{sanchez2022zclaim}。Han 等人提出的 P2C2T 进一步指出，异构货币系统间缺乏同时满足原子性、隐私性与实用性的互联机制，并给出仅要求底层链具备签名验证能力的隐私跨链转移方案~\cite{han2025p2c2t}。综上，无脚本隐私链的跨链诉求可概括为三点：去中心化兑换与出入金、与主流链资产互通而不破坏隐私/可互换性，以及在脚本能力缺失条件下仍能完成原子交换。
然而，多链生态也带来了资产和状态割裂的问题~\cite{JSGG202402003}。某一条链上的账户余额、交易记录和合约状态通常不能被另一条链直接读取或修改，用户如果希望在不同链之间完成资产交换，不能简单地调用一次普通转账交易，而必须借助跨链协议完成多链状态协调~\cite{zamyatin2021sok}。

跨链资产交换的本质并不是把同一个资产实体从一条链物理搬运到另一条链，而是在两个相互独立的账本系统之间建立可信的价值对应关系。以 Alice 与 Bob 的交换为例，Alice 希望在源链上支付一定数量的资产，Bob 在目标链上收到等值资产。由于两条链的交易确认相互独立，如果没有额外机制，Alice 支付后 Bob 可能无法收到资产，或者 Bob 收到资产后 Alice 的支付无法完成。跨链系统必须解决”谁先执行””如何证明对方已执行””失败后如何回滚”等核心问题。

在传统中心化交易平台中，这类问题通常由平台托管资产并撮合交易解决，但这种方式引入了中心化信任风险。平台一旦出现作恶、被攻击、运营失败或提现限制，用户资产可能遭受损失。去中心化跨链交换希望通过链上合约、密码学锁定、事件监听和中继机制降低对单一中心化实体的依赖，使资产释放条件由可验证的协议规则控制~\cite{zamyatin2021sok}。本文所研究的系统正是面向这一问题的原型实现，旨在通过多角色协作和密码学条件锁定，为跨链资产交换提供一种结构清晰、安全可分析的设计方案。

\section{研究意义}

跨链资产交换的研究具有较强的现实意义和工程价值。从需求侧看，多链生态已经成为区块链应用发展的常态~\cite{JSGG202402003,lohachab2021towards}，用户希望将一条链上的原生资产转换为另一条链上的资产，或者在不同链之间使用同一套应用服务。如果缺少跨链能力，各链将形成相互割裂的价值孤岛，限制区块链应用的扩展和用户体验的统一。

从技术方案看，现有跨链方案大致可分为中继链验证、侧链/桥托管和原子交换三类~\cite{Zhu2024,3845}。中继链验证安全性较强但实现成本高，对链的合约能力和验证能力要求较高，Cosmos~\cite{kwon2019cosmos} 和 Polkadot~\cite{wood2016polkadot} 是这一方向的典型代表；侧链/桥托管实现相对简单但引入了中心化信任假设~\cite{back2014enabling,singh2020sidechain}，削弱了区块链系统本身强调的去中心化特性。基于 HTLC 的原子交换~\cite{10.1145/3212734.3212736} 提供了一种较轻量的折中方案，通过 hashLock 和 timeLock 将两侧资产的领取条件绑定在同一个 secret 上，能够在不完全信任对手方的情况下完成交换~\cite{poon2016bitcoin}，同时结合公证人协调机制~\cite{xiong2022notary,9860209} 可以兼顾流程自动化与安全性。

与部分依赖流动性提供者或中继节点在目标链垫付资产的跨链桥不同~\cite{yin2022bool}，本文将 Notary 限定为事件监听和消息协调角色。Notary 不承担资产托管或目标链流动性提供职责，其主要作用是观察源链事件、解析跨链参数并推动目标链创建对应锁定。资金锁定和释放由两侧 Manager/TokenPool 以及 HTLC 或 adaptor signature 条件机制约束。换言之，Notary 在真实链上仍需要承担提交交易所需的 gas 成本，但不需要为用户跨链交换提供目标链资产流动性，这有助于降低中继节点的资金门槛，并使系统安全性更多依赖可验证的合约条件和密码学关系。

从技术难点看，异构链兼容是跨链系统必须面对的挑战。不同链在脚本表达能力上差异巨大：EVM 类链可以直接部署智能合约，通过合约状态记录 hashLock、timeLock、swapId 和资产锁定情况，HTLC 所需的哈希锁判断和条件分支都能在链上显式表达；而 Monero、ZCash 屏蔽交易等无脚本链以纯密码学方式构造交易，不提供通用脚本能力，HTLC 依赖的链上条件逻辑根本无法迁移~\cite{9833731}。相关工作也从应用侧印证了这一缺口：Thyagarajan 等人~\cite{thyagarajan2022universal} 与 Han 等人~\cite{han2025p2c2t} 均将“仅依赖签名验证”作为覆盖隐私币与异构账本的最小能力假设；Thyagarajan 等人~\cite{thyagarajan2022paymo}、Hoenisch 等人~\cite{hoenisch2022lightswap} 与 Sanchez 等人~\cite{sanchez2022zclaim} 的研究则分别表明，门罗币与 Zcash 用户需要在不硬分叉、尽量保持隐私与可互换性的前提下完成与比特币等链的原子交换或资产迁移。如果跨链系统只支持一种合约模型，就无法覆盖这一类只能签名、不能执行脚本的链。本文将 HTLC 与 Schnorr adaptor signature~\cite{9833731} 放在统一跨链架构下讨论：对于合约能力较强的链使用链上 HTLC，对于无脚本链则将条件锁定逻辑放入签名补全过程中，使条件化交易在仅具备签名能力的链上同样能够实现，从而在统一流程下覆盖从智能合约链到无脚本链的不同执行能力。

除了算法理论层面提出跨链资产交换协议之外，本文同样开发了能在主流区块链系统间真实跨链转账的原型系统，具有系统实现方面的重要意义。项目中包含 Solidity 智能合约、Python 后端协调服务、scriptless 场景下的 adaptor signature 模拟模块（以 BTC 作为替代实现）和前端可视化页面，能够较完整地展示跨链资产交换从协议设计到代码实现再到流程演示的全过程。

\section{本文主要工作}

本文围绕跨链资产交换原型系统，从架构设计、核心机制实现和安全性分析三个层面展开工作。

在架构设计层面，本文提出了 Alice - ManagerA - Notary - ManagerB - Bob 的跨链角色模型。该模型将用户、资金池管理者、公证中继节点和目标链接收方明确区分，使系统能够清晰表达资产控制边界与消息协调之间的关系，避免将跨链流程简单等同于单个合约调用。系统将跨链资产交换拆分为用户发起、资金池锁定、事件监听、目标链响应、秘密揭示、资产释放与超时退款等阶段，每个阶段由对应角色驱动。

在核心机制实现层面，本文针对不同链的脚本表达能力设计了两种条件锁定方式。对于支持智能合约的链，系统通过 HTLC 实现资产锁定，使用 hashLock 约束领取条件、timeLock 约束退款条件、swapId 关联同一笔跨链请求，形成源链与目标链的状态配对。对于 Monero、ZCash 屏蔽交易等无脚本链，由于 HTLC 所需的哈希锁判断和条件分支无法在链上表达，系统引入 Schnorr adaptor signature，将完成条件嵌入签名补全过程，以纯密码学方式实现条件化交易，使这类只能签名、不能执行脚本的链也能参与跨链原子交换。在两种机制之上，Notary 作为公证人监听源链合约事件，读取跨链参数并协调目标链 Manager 创建对应锁定，承担跨链消息传递和流程推动的职责。工程实现方面，智能合约模块负责资产池和 HTLC 状态管理，Python 后端模块负责跨链协调，前端模块用于展示跨链流程中各角色状态、密码学参数和交易进度的变化。

在安全性分析层面，本文对系统的原子性、可退款性、Notary 信任边界、Manager 约束和异构兼容能力进行了讨论，说明该系统在原型层面能够满足跨链资产交换的基本安全要求。

本文其余部分组织如下：第二章介绍 HTLC 和 Schnorr adaptor signature 的相关技术背景；第三章阐述系统总体设计，包括角色模型、模块划分和跨链流程；第四章详细描述核心机制的设计与实现；第五章进行系统安全性分析；第六章展示系统实现与实验验证；第七章总结全文并讨论后续工作方向。

